% Part 1, Problem 17: Density of irrationals - Constructive Proof
% Hard

\begin{problem}[Irrationals Between Any Two Rationals]
Prove that for any two distinct rational numbers $p$ and $q$, there exists an irrational number $x$ such that
\[
p < x < q
\]
or
\[
q < x < p.
\]
\end{problem}

\begin{hint}
Attempt the proof independently first. Focus on the key theorem, algebraic transformation, or contradiction setup that links the hypothesis to the target conclusion.
\end{hint}

\begin{solution}
\textbf{Constructive Proof}

Without loss of generality, assume $p<q$.
Define
\[
x=p+\frac{\sqrt{2}}{2}(q-p).
\]

Since $0<\frac{\sqrt{2}}{2}<1$ and $q-p>0$, multiplying by $(q-p)$ gives
\[
0<\frac{\sqrt{2}}{2}(q-p)<q-p.
\]
Adding $p$ throughout,
\[
p<x<q.
\]

It remains to show that $x$ is irrational. Suppose, for contradiction, that $x$ is rational.
Then $x-p$ is rational, so
\[
x-p=\frac{\sqrt{2}}{2}(q-p)
\]
is rational.
Since $q-p$ is a nonzero rational number, dividing by it gives
\[
\frac{\sqrt{2}}{2}=\frac{x-p}{q-p},
\]
which is rational. Hence $\sqrt{2}$ would be rational, a contradiction.

Therefore $x$ is irrational, and we have found an irrational number strictly between $p$ and $q$.

\hfill $\square$
\end{solution}

\begin{takeaways}
\begin{itemize}
\item \textbf{Constructive Method:} The proof gives an explicit irrational number between the two rationals
\item \textbf{Scaling Trick:} Multiply the rational gap $(q-p)$ by an irrational factor between $0$ and $1$
\item \textbf{Density of Irrationals:} Irrational numbers occur between any two distinct rationals
\item \textbf{Companion Result:} Together with density of rationals, this shows the two sets are interwoven on the real line
\item For any two distinct rational numbers $p$ and $q$, there exists infinitely many irrational numbers $x$ such that $p < x < q$ or $q < x < p$.
\item Because the two sets ($\mathbb{Q}$ and $\mathbb{R} \setminus \mathbb{Q}$) of numbers are interwoven, no matter how closely you zoom in on the real line, you will always find both rational and irrational numbers.
\item While they are perfectly interwoven and mutually dense, they are not equal. The rationals are countable, meaning the set of all rational numbers can be put into a one-to-one correspondence with the natural numbers, while the irrationals are uncountable, meaning there are "more" irrationals than rationals in a precise mathematical sense.
\item The probability of randomly selecting a rational number the interval $[0, 1]$ is zero! Learn "Measure Theory" to understand this concept of "size" of sets in a more rigorous way.
\end{itemize}
\end{takeaways}
